\documentclass[12pt]{article}
\usepackage[utf8]{inputenc}
\usepackage[margin=1in]{geometry}
\usepackage{amsmath}
\usepackage{graphicx}
\usepackage{booktabs}
\usepackage{hyperref}
\usepackage[round]{natbib}

\title{Mapping the Core Site for Active Expiration Generation Along the Rostro-Caudal Axis of the Lateral Parafacial Region Using Focal Bicuculline Injections}
\author{Author A$^{1}$, Author B$^{1,2}$, Author C$^{2}$, Author D$^{1}$ \\
\small $^{1}$Department of Physiology, Faculty of Medicine, University of S\~{a}o Paulo, S\~{a}o Paulo, Brazil \\
\small $^{2}$Neuroscience Graduate Program, University of S\~{a}o Paulo, S\~{a}o Paulo, Brazil}
\date{}

\begin{document}
\maketitle

\begin{abstract}
The expiratory oscillator in the lateral parafacial region (pFL) of the ventral medulla is silent at rest and lacks a definitive anatomical marker, making precise localization difficult. Prior pharmacological and optogenetic studies targeted locations near the caudal tip of the facial nucleus (VIIc) spanning roughly $-0.2$ to $+0.5$ mm, but chemogenetic inhibition at these sites only partially suppressed expiratory output. Here, we systematically mapped respiratory responses to bilateral bicuculline methochloride injections (200 $\mu$M, 200 nL/side) at five rostro-caudal positions ($-0.2$ to $+0.8$ mm from VIIc) in urethane-anesthetized adult male Sprague-Dawley rats ($n = 35$). All injection sites recruited abdominal (ABD) activity, but response robustness varied systematically along the axis. The most rostral sites ($+0.6$ and $+0.8$ mm) produced the strongest and longest-lasting responses: tidal volume increased 29\% at $+0.6$ mm ($p = 0.02$), minute ventilation increased 42\% ($p = 0.001$), and the ventilation-to-oxygen-consumption ratio increased 73--82\% ($p = 0.001$). ABD response duration was approximately 17 minutes at rostral sites versus only 2.4 minutes at $-0.2$ mm ($p = 0.043$). The $+0.6$ mm site uniquely initiated responses after the first injection and was the only location producing a post-inspiratory ABD peak (5/6 rats). Multivariate trajectory analysis of the respiratory cycle differentiated injection sites more effectively than any single variable, confirming the rostral pFL as the core region for active expiration generation. cFos analysis confirmed local activation ($\sim$300 $\mu$m spread) without engaging PHOX2B$^+$ RTN neurons above baseline.
\end{abstract}

\section{Introduction}

Breathing rhythm generation depends on brainstem oscillators that produce inspiratory and expiratory motor patterns. While the inspiratory oscillator in the pre-B\"{o}tzinger complex is well characterized, the expiratory oscillator has been more difficult to locate precisely \citep{feldman2013breathing, janczewski2006role}. This oscillator, thought to reside in the lateral parafacial region (pFL) of the ventral medulla, is conditionally active: it is silent during quiet breathing at rest but engages during exercise, hypercapnia, and other conditions requiring active expiration \citep{pagliardini2011active, abdala2009abdominal}.

The conditional silence of the expiratory oscillator, combined with the absence of a definitive anatomical marker for its constituent neurons, has made precise localization challenging. Prior pharmacological and optogenetic studies have targeted locations near the caudal tip of the facial nucleus (VIIc), spanning approximately $-0.2$ to $+0.5$ mm from VIIc \citep{pagliardini2011active, huckstepp2015pfl}. However, chemogenetic inhibition at VIIc only partially suppressed expiratory output, with abdominal recruitment still occurring during sleep \citep{silva2016inhibition}. Similarly, inhibition at $+0.5$ mm reduced but did not eliminate bicuculline/strychnine-induced ABD signals \citep{anderson2016pfl}. These observations suggest that the core of the expiratory oscillator may lie at a different rostro-caudal position than previously assumed.

No systematic study has quantified respiratory responses across the full rostro-caudal extent of the pFL to identify which sub-region drives the most potent active expiration. Moreover, prior analyses typically examined individual respiratory variables in isolation, which may fail to capture the complex, multi-dimensional changes in the breathing pattern produced by pFL activation. Here, we address these gaps by (1) mapping responses to bilateral bicuculline injections at five positions spanning $-0.2$ to $+0.8$ mm from VIIc, (2) measuring a comprehensive set of respiratory variables including phase-specific analyses, and (3) developing a multivariate trajectory analysis that captures the full respiratory cycle as a three-dimensional object to differentiate injection sites.

\section{Methods}

\subsection{Animals and Surgical Preparation}

Adult male Sprague-Dawley rats ($n = 35$; mean weight $340.4 \pm 13.2$ g) were anesthetized with urethane (1.5--1.7 g/kg IV after 5\% isoflurane induction). Animals were instrumented with tracheal cannulae for airflow measurement, EMG electrodes in oblique abdominal and diaphragm muscles, and femoral vein cannulae. Supplemental O$_2$ (30\%) was provided and body temperature maintained at $37 \pm 1^\circ$C. Bilateral mid-cervical vagotomy with 2 mm resection was performed.

\subsection{Injection Protocol}

Bicuculline methochloride (200 $\mu$M in HEPES with fluorobeads) or HEPES vehicle was pressure-injected bilaterally (200 nL/side, 100 nL/min, 30 $\mu$m pipette tip) into the pFL at five rostro-caudal positions (Table~\ref{tab:groups}).

\begin{table}[htbp]
\centering
\caption{Injection group allocation.}
\label{tab:groups}
\begin{tabular}{lcc}
\toprule
Position (mm from VIIc) & $n$ (bicuculline) & $n$ (HEPES control) \\
\midrule
$-0.2$ & 5 & -- \\
$+0.1$ & 7 & -- \\
$+0.4$ & 5 & -- \\
$+0.6$ & 6 & -- \\
$+0.8$ & 5 & -- \\
CTRL & -- & 7 \\
Total & 28 & 7 \\
\bottomrule
\end{tabular}
\end{table}

\subsection{Respiratory Measurements}

Recorded signals included tracheal airflow, integrated ABD EMG, integrated diaphragm (DIA) EMG, and end-tidal CO$_2$/O$_2$. Derived metrics included respiratory rate, tidal volume (VT, mL/kg), minute ventilation (VE, mL/min/kg), O$_2$ consumption (VO$_2$, calculated per \citealt{depocas1957metabolic}), and VE/VO$_2$ ratio. Phase-specific analysis computed normalized area under mean-cycle signals for late-E, inspiratory, and post-inspiratory phases.

\subsection{Multivariate Trajectory Analysis}

Mean respiratory cycles (airflow, DIA EMG, ABD EMG) were projected into 3D space at each 2-minute time bin. Procrustes-like distance metrics measured trajectory deviations from baseline, enabling simultaneous assessment of all respiratory signals across the full breathing cycle.

\subsection{Histology}

Injection sites were verified with fluorobeads and immunohistochemistry for cFos, PHOX2B, ChAT, and TH \citep{stornetta2006phox2b, guyenet2019rtn}.

\subsection{Statistics}

One-way ANOVA with Tukey or Bonferroni post hoc, repeated-measures ANOVA, Kruskal-Wallis with Dunn post hoc (Sidak correction), and Wilcoxon rank-sum tests were used. Effect sizes are reported as $\eta^2$.

\section{Results}

\subsection{All pFL Sites Recruit ABD Activity, but Robustness Varies}

Bicuculline injections at all five rostro-caudal positions recruited late-expiratory (late-E) ABD activity, confirming the broad extent of the pFL. However, the robustness and qualitative features of the response varied systematically along the axis (Table~\ref{tab:abd}).

\begin{table}[htbp]
\centering
\caption{ABD response characteristics by injection site.}
\label{tab:abd}
\begin{tabular}{lcccc}
\toprule
Position & Responders & Duration (min) & Tonic component & Post-I ABD \\
\midrule
$-0.2$ mm & 3/5 & $2.4 \pm 1.1$ & No & No \\
$+0.1$ mm & 7/7 & -- & No & No \\
$+0.4$ mm & 5/5 & -- & No & No \\
$+0.6$ mm & 6/6 & $17.6 \pm 2.7$ & Yes & Yes (5/6) \\
$+0.8$ mm & 5/5 & $17.1 \pm 3.3$ & Yes & No \\
\bottomrule
\end{tabular}
\end{table}

ABD response duration differed significantly among groups (one-way ANOVA: $p = 0.043$, $\eta^2 = 0.41$), with rostral sites ($+0.6$, $+0.8$ mm) showing approximately seven-fold longer responses than the most caudal site ($-0.2$ mm; Tukey post hoc: $p < 0.05$). The $+0.6$ mm site was unique in two respects: it initiated responses after the first injection (before the second), and 5 of 6 rats showed a post-inspiratory ABD peak not observed at other sites.

\subsection{Rostral Sites Produce Strongest Ventilatory Effects}

Standard respiratory variables showed systematic rostro-caudal gradients. Tidal volume peaked at $+0.6$ mm with a 29\% increase from baseline (one-way ANOVA: $p = 0.02$, $\eta^2 = 0.42$; Bonferroni: $-0.2$ vs. $+0.6$ mm: $p = 0.002$). Minute ventilation increased 42\% at $+0.6$ mm (repeated-measures ANOVA: $p = 0.001$, $\eta^2 = 0.59$), driven by increased VT overcoming the drop in respiratory rate (Table~\ref{tab:ventilation}).

\begin{table}[htbp]
\centering
\caption{Summary of ventilatory effect sizes across measures.}
\label{tab:ventilation}
\begin{tabular}{lcccl}
\toprule
Measure & ANOVA $p$ & $\eta^2$ & Best site & \% Change \\
\midrule
ABD duration & 0.043 & 0.41 & $+0.6$/$+0.8$ & 7$\times$ vs. $-0.2$ \\
Max tidal volume & 0.02 & 0.42 & $+0.6$ & $+29\%$ \\
Max minute ventilation & 0.001 & 0.59 & $+0.6$ & $+42\%$ \\
Min O$_2$ consumption & 0.0001 & 0.59 & $+0.6$/$+0.8$ & $-33\%$ \\
Max VE/VO$_2$ & 0.001 & 0.55 & $+0.8$ & $+82\%$ \\
\bottomrule
\end{tabular}
\end{table}

O$_2$ consumption decreased selectively at rostral sites ($-33\%$ at $+0.6$ mm; repeated-measures ANOVA: $p = 0.0001$, $\eta^2 = 0.59$), while caudal sites showed minimal change. The combination of increased VE and decreased VO$_2$ produced hyperventilation at rostral sites, with VE/VO$_2$ increasing 73\% at $+0.6$ mm and 82\% at $+0.8$ mm ($p = 0.001$, $\eta^2 = 0.55$).

Respiratory rate decreased across all groups, driven by prolonged expiration (TE: baseline $1.029 \pm 0.161$ s, response $1.313 \pm 0.188$ s; Wilcoxon $Z = 4.49$, $p = 0.001$) and shortened inspiration (Ti: baseline $0.318 \pm 0.043$ s, response $0.279 \pm 0.034$ s; $Z = 3.24$, $p = 0.001$). The largest rate decrease occurred at rostral sites ($-25\%$ at $+0.6$ mm).

\subsection{Phase-Specific Analysis}

Late-E phase analysis confirmed that injections at all locations induced negative airflow deflections (active expiration). The $+0.6$ mm site showed significantly elevated late-E airflow versus control from 4 to 12 minutes post-injection (Kruskal-Wallis $H(5) = 14.66$--$18.35$, $p < 0.02$; Dunn $p < 0.001$). Inspiratory phase airflow was also elevated at $+0.6$ mm versus control (Kruskal-Wallis $H(4) = 13.25$, $p = 0.021$). Post-inspiratory ABD activity was most pronounced at $+0.6$ and $+0.8$ mm groups.

\subsection{Multivariate Trajectory Analysis Differentiates Injection Sites}

Three-dimensional respiratory trajectories (airflow $\times$ DIA EMG $\times$ ABD EMG) showed that rostral injection sites produced the largest and most sustained deviations from baseline. The combined trajectory distance metric differentiated injection sites more effectively than any single respiratory variable, with significant between-group differences at multiple time points (Kruskal-Wallis $p < 0.05$). Control injections showed minimal trajectory deviations.

\subsection{cFos Confirms Local Activation Without RTN Engagement}

Bicuculline injections produced elevated cFos$^+$ cell counts at injection sites ($104.5 \pm 5.1$ cells/hemisection) compared to HEPES controls ($44.7 \pm 4.0$), with counts dropping to baseline ($44.8 \pm 1.2$) beyond injection boundaries. Fluorobead spread was approximately 300 $\mu$m. Critically, cFos$^+$/PHOX2B$^+$ double-labeled cells did not differ between bicuculline and control groups ($2.4 \pm 0.4$ vs. $2.3 \pm 1.0$ per hemisection), confirming that RTN chemosensory neurons were not activated above baseline \citep{stornetta2006phox2b, guyenet2019rtn}. No overlap was observed with TH$^+$ (C1/A5) or ChAT$^+$ neurons.

\section{Discussion}

This study provides the first systematic rostro-caudal mapping of active expiration responses in the lateral parafacial region, revealing that the most rostral sites ($+0.6$ and $+0.8$ mm from VIIc) produce the strongest, longest-lasting, and most physiologically comprehensive respiratory changes.

\subsection{Rostral pFL as the Core Expiratory Oscillator Site}

Multiple lines of evidence converge on the rostral pFL as the core site for active expiration generation \citep{feldman2013breathing, janczewski2006role}. The $+0.6$ mm site was distinguished by: (1) the longest ABD response duration ($\sim$17 minutes vs. $\sim$2.4 minutes at $-0.2$ mm), (2) the largest increases in tidal volume and minute ventilation, (3) unique initiation of responses after the first injection, and (4) being the only site producing a post-inspiratory ABD peak. These features suggest that the $+0.6$ mm location contains the densest concentration of expiratory oscillator neurons within the pFL \citep{pagliardini2011active, huckstepp2015pfl}.

This rostral localization is consistent with partial suppression results from prior chemogenetic inhibition studies at more caudal locations \citep{silva2016inhibition, anderson2016pfl}. If the core oscillator lies at $+0.6$ mm, inhibition at VIIc ($0$ mm) or $+0.5$ mm would only partially overlap with the critical neuronal population, explaining incomplete suppression.

\subsection{Physiological Significance of Rostral pFL Activation}

The combination of increased minute ventilation and decreased oxygen consumption at rostral sites produced hyperventilation (VE/VO$_2$ increase of 73--82\%), a response absent at caudal sites. This suggests that the rostral pFL drives a coordinated ventilatory response that goes beyond simply adding an expiratory phase to the breathing cycle---it reconfigures the entire pattern of respiratory muscle activation and metabolic coupling \citep{abdala2009abdominal, deruyver2017pfl}.

\subsection{Multivariate Analysis as a Tool for Respiratory Phenotyping}

The 3D trajectory analysis proved more effective than any single variable for differentiating injection sites, because it captures simultaneous changes across multiple signals and respiratory phases. This approach could be valuable for characterizing respiratory phenotypes in other experimental contexts, such as chemogenetic or optogenetic manipulations of respiratory circuits.

\subsection{Limitations}

The study was conducted under urethane anesthesia with bilateral vagotomy, which alters respiratory pattern and eliminates vagal feedback. The bicuculline concentration (200 $\mu$M) and volume (200 nL) produce local but not perfectly focal activation. Individual late-E peak amplitude and area comparisons did not reach statistical significance despite clear trends, likely reflecting power limitations with group sizes of 5--7.

\section{Conclusion}

Systematic rostro-caudal mapping of the lateral parafacial region identifies the site at $+0.6$ mm from the caudal tip of the facial nucleus as the core location for active expiration generation, producing the strongest, longest-lasting, and most physiologically comprehensive respiratory responses to GABAergic disinhibition. Multivariate respiratory cycle analysis provides the most sensitive tool for distinguishing the respiratory effects of different pFL sub-regions. These findings refine the anatomical localization of the expiratory oscillator and have implications for targeting this region in future mechanistic and therapeutic studies.

\bibliographystyle{plainnat}
\bibliography{references}

\end{document}
